\documentclass[11pt,a4paper]{article}
\usepackage[utf8]{inputenc}
\usepackage[T1]{fontenc}
\usepackage[spanish,es-nodecimaldot]{babel}
\usepackage{lmodern,microtype}
\usepackage[a4paper,landscape,margin=1.2cm]{geometry}
\usepackage{longtable,booktabs,array,xcolor}
\usepackage[unicode,colorlinks=true,urlcolor=blue!45!black]{hyperref}
\pagestyle{empty}
\setlength{\parindent}{0pt}
\setlength{\tabcolsep}{4.5pt}
\setlength{\LTpre}{0pt}
\setlength{\LTpost}{0pt}
\setlength{\emergencystretch}{2em}
\renewcommand{\arraystretch}{1.10}
\newcolumntype{L}[1]{>{\raggedright\arraybackslash}p{#1}}
\makeatletter
\newcommand{\sourceurl}[2]{\hyper@linkurl{#2}{#1}}
\makeatother
\begingroup\catcode`\%=12
\expandafter\gdef\csname srcS01\endcsname{\sourceurl{file:///C:/Users/orlan/Desktop/SanBartolo/global_economics/book2026-2027/seminars/term_1/second_learning_evidence/C5%2C6%2C7%2C8%2C9%2C10.%20Nicol%C3%A1s%20%C3%81lvarez%2C%20Thomas%20%C3%81lvarez%2C%20Juan%20Felipe%20Ojeda..pdf}{Álvarez y Ojeda, C5--C10}}
\endgroup
\begingroup\catcode`\%=12
\expandafter\gdef\csname srcS02\endcsname{\sourceurl{file:///C:/Users/orlan/Desktop/SanBartolo/global_economics/book2026-2027/seminars/term_1/second_learning_evidence/Criterio_10.pdf}{Ayala y Galán, acumulativo C3--C10}}
\endgroup
\begingroup\catcode`\%=12
\expandafter\gdef\csname srcS03\endcsname{\sourceurl{file:///C:/Users/orlan/Desktop/SanBartolo/global_economics/book2026-2027/seminars/term_1/second_learning_evidence/Encuesta_SBLM%20%282%29.pdf}{Ayala y Galán, Encuesta SBLM}}
\endgroup
\begingroup\catcode`\%=12
\expandafter\gdef\csname srcS04\endcsname{\sourceurl{file:///C:/Users/orlan/Desktop/SanBartolo/global_economics/book2026-2027/seminars/term_1/second_learning_evidence/c5c6c7c8.pdf}{González y Barón, C5--C6}}
\endgroup
\begingroup\catcode`\%=12
\expandafter\gdef\csname srcS05\endcsname{\sourceurl{file:///C:/Users/orlan/Desktop/SanBartolo/global_economics/book2026-2027/seminars/term_1/second_learning_evidence/Proyecto_C7-C8.pdf}{González y Barón, Proyecto C7--C8}}
\endgroup
\begingroup\catcode`\%=12
\expandafter\gdef\csname srcS06\endcsname{\sourceurl{file:///C:/Users/orlan/Desktop/SanBartolo/global_economics/book2026-2027/seminars/term_1/second_learning_evidence/Lucio_IsaC_Criterio_5-10.pdf}{Isa y Lucio, hábitos de sueño}}
\endgroup
\begingroup\catcode`\%=12
\expandafter\gdef\csname srcS07\endcsname{\sourceurl{file:///C:/Users/orlan/Desktop/SanBartolo/global_economics/book2026-2027/seminars/term_1/second_learning_evidence/SeminariosMiguel_Ignacio_tex%20%283%29.pdf}{Caballero y Rodríguez, proyecto}}
\endgroup
\begingroup\catcode`\%=12
\expandafter\gdef\csname srcS08\endcsname{\sourceurl{file:///C:/Users/orlan/Desktop/SanBartolo/global_economics/book2026-2027/seminars/term_1/second_learning_evidence/C1-C10%2C%20Sebastian%20Lozano%2C%20Alejandro%20Camargo.pdf}{Lozano y Camargo, C1--C10}}
\endgroup
\begingroup\catcode`\%=12
\expandafter\gdef\csname srcS09\endcsname{\sourceurl{file:///C:/Users/orlan/Desktop/SanBartolo/global_economics/book2026-2027/seminars/term_1/second_learning_evidence/Segunda_evidencia_de_apredizaje__C6_C10_.pdf}{Del Cairo y Clavijo, residuos}}
\endgroup
\begingroup\catcode`\%=12
\expandafter\gdef\csname srcS10\endcsname{\sourceurl{file:///C:/Users/orlan/Desktop/SanBartolo/global_economics/book2026-2027/seminars/term_1/second_learning_evidence/Di%CC%81az%2C%20Maldonado%20%26%20Lo%CC%81pez%20C5%20%26%20C6.pdf}{Díaz, Maldonado y López, C5--C6}}
\endgroup
\begingroup\catcode`\%=12
\expandafter\gdef\csname srcS11\endcsname{\sourceurl{file:///C:/Users/orlan/Desktop/SanBartolo/global_economics/book2026-2027/seminars/term_1/second_learning_evidence/C3%2C%20Andre%CC%81s%20Maldonado%2C%20Carlos%20Lo%CC%81pez%2C%20Thomas%20Di%CC%81az.pdf}{Encuesta de graderías, 30 estudiantes}}
\endgroup
\begingroup\catcode`\%=12
\expandafter\gdef\csname srcS12\endcsname{\sourceurl{file:///C:/Users/orlan/Desktop/SanBartolo/global_economics/book2026-2027/seminars/term_1/second_learning_evidence/Lo%CC%81pez%2C%20Di%CC%81az%20%26%20Maldonado%20C7%20%26%20C8.docx}{López, Díaz y Maldonado, C7--C8 (DOCX)}}
\endgroup
\begingroup\catcode`\%=12
\expandafter\gdef\csname srcS13\endcsname{\sourceurl{file:///C:/Users/orlan/Desktop/SanBartolo/global_economics/book2026-2027/seminars/term_1/second_learning_evidence/Maldonado%2C%20Lopez%2C%20Diaz%20C9%20%26%20C10%20.docx}{Maldonado, López y Díaz, C9--C10 (DOCX)}}
\endgroup
\begingroup\catcode`\%=12
\expandafter\gdef\csname srcS14\endcsname{\sourceurl{file:///C:/Users/orlan/Desktop/SanBartolo/global_economics/book2026-2027/seminars/term_1/second_learning_evidence/Seminarios%20%281%29.pdf}{Guampe y Ospina, cafetería}}
\endgroup
\begingroup\catcode`\%=12
\expandafter\gdef\csname srcS15\endcsname{\sourceurl{file:///C:/Users/orlan/Desktop/SanBartolo/global_economics/book2026-2027/seminars/term_1/second_learning_evidence/Seminarios_2026_2027%20%282%29.pdf}{Julián Maldonado, organización académica}}
\endgroup
\begingroup\catcode`\%=12
\expandafter\gdef\csname srcS16\endcsname{\sourceurl{file:///C:/Users/orlan/Desktop/SanBartolo/global_economics/book2026-2027/seminars/term_1/second_learning_evidence/C5.pdf}{Ospina y Méndez, C5--C6}}
\endgroup
\begingroup\catcode`\%=12
\expandafter\gdef\csname srcS17\endcsname{\sourceurl{file:///C:/Users/orlan/Desktop/SanBartolo/global_economics/book2026-2027/seminars/term_1/second_learning_evidence/C7..pdf}{Ospina y Méndez, C7--C8}}
\endgroup
\begingroup\catcode`\%=12
\expandafter\gdef\csname srcS18\endcsname{\sourceurl{file:///C:/Users/orlan/Desktop/SanBartolo/global_economics/book2026-2027/seminars/term_1/second_learning_evidence/C9.pdf}{Ospina y Méndez, C9--C10}}
\endgroup
\begingroup\catcode`\%=12
\expandafter\gdef\csname srcS19\endcsname{\sourceurl{file:///C:/Users/orlan/Desktop/SanBartolo/global_economics/book2026-2027/seminars/term_1/second_learning_evidence/pedro_y_tape.pdf}{Pedro y Santiago, baños}}
\endgroup
\begingroup\catcode`\%=12
\expandafter\gdef\csname srcS20\endcsname{\sourceurl{file:///C:/Users/orlan/Desktop/SanBartolo/global_economics/book2026-2027/seminars/term_1/second_learning_evidence/Seminarios%20Tom%C3%A1s%20Rinc%C3%B3n.pdf}{Tomás Rincón, filas}}
\endgroup
\makeatletter
\newcommand{\srclink}[1]{\begingroup\renewcommand{\sourceurl}[2]{\hyper@linkurl{#1}{##1}}\csname src#1\endcsname\endgroup}
\makeatother
\begin{document}
\begin{longtable}{@{}L{1.25cm}L{5.1cm}L{7.6cm}L{4cm}L{7.9cm}@{}}
\toprule
\textbf{Group} & \textbf{Members} & \textbf{Work C5--C10} & \textbf{Problem} & \textbf{Solutions} \\
\midrule
\endfirsthead
\toprule
\textbf{Group} & \textbf{Members} & \textbf{Work C5--C10} & \textbf{Problem} & \textbf{Solutions} \\
\midrule
\endhead
\bottomrule
\endfoot
\bottomrule
\endlastfoot
1 & Thomas Álvarez; Nicolás Álvarez Ospina; Juan Felipe Ojeda Hurtado. Thomas no figura en el CSV. & C5--C6: recorrido y priorización; C7--C8: turnos equitativos, acción colectiva y señalética; C9--C10: once ideas y contraste de supuestos; autoría y escucha anónimas; confusión espacial pendiente. & Colados y desorden al comprar alimentos o usar microondas reducen el descanso. & Proponen delimitar filas con supervisión estudiantil, no atender a colados y separar los microondas con fichas de tiempo. {\footnotesize \srclink{S01}} \\
\midrule
2 & Patrik Ayala Palenik (variante Patrick); Gabriel Hernando Galán Osorio. & C5--C6: testimonios de espera y necesidad prioritaria; matriz comparativa no identificada en las fuentes; C7--C8: flujos y convivencia; C9--C10: cinco propuestas y contraste de acceso sin teléfono. & Espera y conflictos por colados al comprar o reclamar en Asobartolina. & Proponen una fila exprés para pedidos prepagados, turnos por colores y una campaña gráfica de respeto de la fila. {\footnotesize \srclink{S02}, \srclink{S03}} \\
\midrule
3 & Sebastián González Mojica; Emilio José Barón Romero. & C5--C6: cuatro hallazgos y priorización de claridad normativa; C7--C8: reformulación, autorregulación y convivencia; C9--C10: registros no identificados en las fuentes; verificaciones institucionales pendientes. & Gestionar la concentración en trabajo autónomo, considerando audífonos, normas y comunicación docente; beneficios de la música no universales. & No se identificaron propuestas desarrolladas en las fuentes revisadas; autorregulación y acuerdos aparecen como conocimientos de C8, sin evidencia de ideación. {\footnotesize \srclink{S04}, \srclink{S05}} \\
\midrule
4 & «Isa» y «Lucio». El CSV vincula Ana Isabella Colorado Parrado y Lucio Barrero Sánchez; equivalencia con esos nombres pendiente. & C5--C6: diez respuestas y comparación de horarios, puntajes parcialmente justificados; C7--C8: hábitos, tiempo y psicología; C9--C10: seis ideas, combinaciones y supuestos; tablas parcialmente legibles; beneficios de estabilidad pendientes. & Horarios de sueño irregulares, tareas y pantallas se relacionan con cansancio escolar. & Proponen un horario personal con cálculo del tiempo de tareas y apagado de pantallas, y un planeador de obligaciones con relajación al terminar. {\footnotesize \srclink{S06}} \\
\midrule
5 & Miguel Ángel Caballero Hernández; Juan Ignacio Rodríguez Rodríguez. & C5--C6: quince encuestados y priorización de clasificación, sin resolver desapariciones; C7--C8: búsqueda, significantes y 5S; C9--C10: seis propuestas y contraste de mantenimiento; celda parcialmente superpuesta; disposición estudiantil pendiente. & La mesa de objetos perdidos se desorganiza y dificulta encontrar pertenencias. & Proponen un inventario virtual en la página del colegio y estudiantes responsables por rotación para ayudar a organizar la mesa. {\footnotesize \srclink{S07}} \\
\noalign{\penalty-10000}
6 & Sebastián Lozano Barreneche; Alejandro Camargo Roldan. & C5--C6: cinco hallazgos y priorización provisional, alternativas sin datos; C7--C8: reformulación, incertidumbre y Kant; C9--C10: seis ideas, combinaciones y ajustes de acceso según medios de pago. & Filas y colados dificultan comer durante el descanso y llegar a clase a tiempo. & Proponen autopedido con carnet, tiquete o efectivo y recogida separada, aún por verificar; pedidos anticipados con turno, conservando pedidos en fila. {\footnotesize \srclink{S08}} \\
\midrule
7 & Pablo José Del Cairo Bolaños; Nicolás Clavijo González. & C5--C6: diez encuestas, siete hallazgos y desempate por viabilidad; C7--C8: clasificación y significantes, rótulos por verificar; C9--C10: seis ideas, aportes cruzados y contraste de comprensión visual y acceso. & Residuos mezclados, basura en el piso y dudas de clasificación aumentan el trabajo de aseo. & Proponen imágenes y colores comunes desde el empaque hasta la caneca, y guías estudiantiles rotativos con comentarios semanales del PAE. {\footnotesize \srclink{S09}} \\
\midrule
8 & Andrés Camilo Maldonado Lemus; Carlos Samuel López Castrillón; Thomas Enrique Díaz Garzón. & C5--C6: encuesta de treinta estudiantes, cuatro hallazgos y priorización del acceso a recipientes; mantenimiento como límite; C7--C10: DOCX accesibles, pendientes de verificación paginada en PDF; anexo de encuesta sin autores. & Acumulación recurrente de envolturas y envases en las graderías, relacionada con hábitos y acceso a canecas. & No se identificó un conjunto de propuestas en los PDF revisados; S10 todavía no diseña solución. Posibles ideas del DOCX S13 pendientes de verificación en PDF. {\footnotesize \srclink{S10}, \srclink{S11}, \srclink{S12}, \srclink{S13}} \\
\midrule
9 & Sofía Guampe Granados (portada: Sofía Guampe G.); Nicolás Ospina Ramírez (portada: Nicolás Ospina R.). & C5--C6: quince encuestas, fechas y mediciones faltantes, priorización provisional sin viabilidad; C7--C8: muestra y escucha, consulta al personal pendiente; C9--C10: seis ideas y contraste; participación propuesta, no verificada como realizada. & Las esperas dificultan comprar, comer y regresar a clase a tiempo. & Proponen explorar pedidos anticipados con reclamo diferenciado e información clara, sin exigir aplicación ni salida anticipada; turnos por llegada con acuerdos y consulta a usuarios. {\footnotesize \srclink{S14}} \\
\noalign{\penalty-10000}
10 & Julián Maldonado, único autor. CSV: «MALDONADO BAQUERO JULIÁN ALEJAND...». Identidad completa pendiente; no se amplía el nombre truncado. & C5--C6: perspectivas de estudiante y profesor, priorización del registro; C7--C8: memoria, atención y UX/UI; C9--C10: revisión atribuida a ChatGPT y contraste de consulta, actualización y acceso; sin colaboración estudiantil documentada. & Olvidos de trabajos, fechas y materiales escolares, con préstamos de compañeros como respuesta temporal. & Propone un mini tablero visual de pendientes en el colegio y un chatbot de WhatsApp con recordatorios y enlaces a una web de tareas, materiales y fechas. {\footnotesize \srclink{S15}} \\
\midrule
11 & Santiago Ospina Ramírez; Juan José Méndez (CSV: MENDEZ RAMOS JUAN JOSE). & C5--C6: diez observados en una fecha, sin generalización, acciones individuales ni citas textuales; priorización provisional; C7--C8: precio y nutrición, materiales/datos faltantes; C9--C10: seis ideas y contraste; algunas exceden D1. & Dificultades para acceder a comida saludable y asequible en la cafetería durante los descansos. & Proponen ampliar el menú saludable, variado y asequible, y alternativas menos procesadas con campañas sobre agua y verduras; preferencia específica y aceptación no demostradas. {\footnotesize \srclink{S16}, \srclink{S17}, \srclink{S18}} \\
\midrule
12 & «Pedro y Santiago». El CSV vincula Pedro Martín Ramírez Barrero y Santiago Tapias López; identidades completas pendientes, especialmente el apellido de Santiago. & C5--C6: perfiles, testimonios y necesidades; matriz comparativa no identificada en las fuentes; C7--C8: reformulación, higiene y responsabilidad colectiva; C9--C10: registros no identificados en las fuentes. & Malos olores, suciedad y daños generan incomodidad y evitan que algunos estudiantes usen los baños. & Plantean mejorar el control y los horarios de aseo, y reparar sanitarios, lavamanos y otras instalaciones deterioradas. {\footnotesize \srclink{S19}} \\
\midrule
13 & Tomás Rincón (CSV: RINCON CASTAÑEDA TOMAS MAURICIO); trabajo individual. & C5--C6: tres hallazgos y comparación con límites de permisos; C7--C8: reformulación, ciudadanía y economía; C9--C10: ideas y contraste individuales, pese a etiqueta grupal; disposición de monitores y capacidad operativa pendientes. & Colados y congestión en compra y reclamo saturan al personal y consumen el descanso. & Propone un canal único demarcado con cintas móviles y monitor; fila rápida para preempacados con efectivo exacto o ficha física, separada de almuerzos. {\footnotesize \srclink{S20}} \\

\end{longtable}
\end{document}
